% Konstruksi HMAC: hashing dua lapis dengan kunci yang dipad.
% Dirujuk dari section-03.tex dengan % Konstruksi HMAC: hashing dua lapis dengan kunci yang dipad.
% Dirujuk dari section-03.tex dengan % Konstruksi HMAC: hashing dua lapis dengan kunci yang dipad.
% Dirujuk dari section-03.tex dengan % Konstruksi HMAC: hashing dua lapis dengan kunci yang dipad.
% Dirujuk dari section-03.tex dengan \input{figures/skema-hmac}.
\begin{figure}[htbp]
    \centering
    \begin{tikzpicture}[
        kotak/.style={draw, rounded corners, align=center, font=\small, inner sep=5pt,
                      minimum width=5.6cm, minimum height=0.75cm},
        hash/.style={kotak, minimum width=2.2cm, fill=orange!12},
        panah/.style={-{Latex[length=2.4mm]}, thick},
        node distance=0.45cm,
    ]
        \node[kotak, fill=blue!6] (a) {$M \;\|\; (K \oplus \text{ipad})$};
        \node[hash, below=of a] (h1) {$H$};
        \node[kotak, fill=gray!8, below=of h1] (b) {$h_1 \;\|\; (K \oplus \text{opad})$};
        \node[hash, below=of b] (h2) {$H$};
        \node[kotak, fill=green!10, below=of h2] (tag) {Tag MAC};

        \draw[panah] (a) -- (h1);
        \draw[panah] (h1) -- (b);
        \draw[panah] (b) -- (h2);
        \draw[panah] (h2) -- (tag);

        \node[right=0.25cm of h1, font=\scriptsize, align=left]
            {$h_1 = H\bigl((K \oplus \text{ipad}) \,\|\, M\bigr)$};
        \node[right=0.25cm of h2, font=\scriptsize, align=left]
            {$\text{Tag} = H\bigl((K \oplus \text{opad}) \,\|\, h_1\bigr)$};
    \end{tikzpicture}
    \caption{Konstruksi HMAC: hashing dua lapis dengan kunci yang dipad memakai ipad dan opad}
    \label{fig:skema-hmac}
\end{figure}
.
\begin{figure}[htbp]
    \centering
    \begin{tikzpicture}[
        kotak/.style={draw, rounded corners, align=center, font=\small, inner sep=5pt,
                      minimum width=5.6cm, minimum height=0.75cm},
        hash/.style={kotak, minimum width=2.2cm, fill=orange!12},
        panah/.style={-{Latex[length=2.4mm]}, thick},
        node distance=0.45cm,
    ]
        \node[kotak, fill=blue!6] (a) {$M \;\|\; (K \oplus \text{ipad})$};
        \node[hash, below=of a] (h1) {$H$};
        \node[kotak, fill=gray!8, below=of h1] (b) {$h_1 \;\|\; (K \oplus \text{opad})$};
        \node[hash, below=of b] (h2) {$H$};
        \node[kotak, fill=green!10, below=of h2] (tag) {Tag MAC};

        \draw[panah] (a) -- (h1);
        \draw[panah] (h1) -- (b);
        \draw[panah] (b) -- (h2);
        \draw[panah] (h2) -- (tag);

        \node[right=0.25cm of h1, font=\scriptsize, align=left]
            {$h_1 = H\bigl((K \oplus \text{ipad}) \,\|\, M\bigr)$};
        \node[right=0.25cm of h2, font=\scriptsize, align=left]
            {$\text{Tag} = H\bigl((K \oplus \text{opad}) \,\|\, h_1\bigr)$};
    \end{tikzpicture}
    \caption{Konstruksi HMAC: hashing dua lapis dengan kunci yang dipad memakai ipad dan opad}
    \label{fig:skema-hmac}
\end{figure}
.
\begin{figure}[htbp]
    \centering
    \begin{tikzpicture}[
        kotak/.style={draw, rounded corners, align=center, font=\small, inner sep=5pt,
                      minimum width=5.6cm, minimum height=0.75cm},
        hash/.style={kotak, minimum width=2.2cm, fill=orange!12},
        panah/.style={-{Latex[length=2.4mm]}, thick},
        node distance=0.45cm,
    ]
        \node[kotak, fill=blue!6] (a) {$M \;\|\; (K \oplus \text{ipad})$};
        \node[hash, below=of a] (h1) {$H$};
        \node[kotak, fill=gray!8, below=of h1] (b) {$h_1 \;\|\; (K \oplus \text{opad})$};
        \node[hash, below=of b] (h2) {$H$};
        \node[kotak, fill=green!10, below=of h2] (tag) {Tag MAC};

        \draw[panah] (a) -- (h1);
        \draw[panah] (h1) -- (b);
        \draw[panah] (b) -- (h2);
        \draw[panah] (h2) -- (tag);

        \node[right=0.25cm of h1, font=\scriptsize, align=left]
            {$h_1 = H\bigl((K \oplus \text{ipad}) \,\|\, M\bigr)$};
        \node[right=0.25cm of h2, font=\scriptsize, align=left]
            {$\text{Tag} = H\bigl((K \oplus \text{opad}) \,\|\, h_1\bigr)$};
    \end{tikzpicture}
    \caption{Konstruksi HMAC: hashing dua lapis dengan kunci yang dipad memakai ipad dan opad}
    \label{fig:skema-hmac}
\end{figure}
.
\begin{figure}[htbp]
    \centering
    \begin{tikzpicture}[
        kotak/.style={draw, rounded corners, align=center, font=\small, inner sep=5pt,
                      minimum width=5.6cm, minimum height=0.75cm},
        hash/.style={kotak, minimum width=2.2cm, fill=orange!12},
        panah/.style={-{Latex[length=2.4mm]}, thick},
        node distance=0.45cm,
    ]
        \node[kotak, fill=blue!6] (a) {$M \;\|\; (K \oplus \text{ipad})$};
        \node[hash, below=of a] (h1) {$H$};
        \node[kotak, fill=gray!8, below=of h1] (b) {$h_1 \;\|\; (K \oplus \text{opad})$};
        \node[hash, below=of b] (h2) {$H$};
        \node[kotak, fill=green!10, below=of h2] (tag) {Tag MAC};

        \draw[panah] (a) -- (h1);
        \draw[panah] (h1) -- (b);
        \draw[panah] (b) -- (h2);
        \draw[panah] (h2) -- (tag);

        \node[right=0.25cm of h1, font=\scriptsize, align=left]
            {$h_1 = H\bigl((K \oplus \text{ipad}) \,\|\, M\bigr)$};
        \node[right=0.25cm of h2, font=\scriptsize, align=left]
            {$\text{Tag} = H\bigl((K \oplus \text{opad}) \,\|\, h_1\bigr)$};
    \end{tikzpicture}
    \caption{Konstruksi HMAC: hashing dua lapis dengan kunci yang dipad memakai ipad dan opad}
    \label{fig:skema-hmac}
\end{figure}
